\documentclass[11pt,a4paper]{article}
\newcommand{\fichetype}{Exercices}
% Thème A4 commun, librement modifiable.
\usepackage[T1]{fontenc}
\usepackage[utf8]{inputenc}
\usepackage[french]{babel}
\usepackage{lmodern,amsmath,amssymb,mathtools,tikz,pgfplots}
\usepackage[margin=2.2cm,headheight=15pt]{geometry}
\usepackage{xcolor,fancyhdr,enumitem,hyperref,needspace}
\usetikzlibrary{arrows.meta,positioning,calc}
\pgfplotsset{compat=1.18,width=11cm,height=6cm}
\definecolor{Ink}{HTML}{202124}
\definecolor{Blue}{HTML}{356B9E}
\definecolor{Muted}{HTML}{737579}
\color{Ink}
\renewcommand{\familydefault}{\sfdefault}
\setlength{\parindent}{0pt}
\setlength{\parskip}{6pt}
\setlist{itemsep=3pt,topsep=4pt}
\pagestyle{fancy}\fancyhf{}
\fancyhead[L]{\small\color{Muted}MATHS LYCÉE}
\fancyhead[R]{\small\color{Muted}2026--2027}
\fancyfoot[L]{\small\color{Muted}\fichetype}
\fancyfoot[R]{\small\thepage}
\renewcommand{\headrulewidth}{0.3pt}
\hypersetup{colorlinks=true,urlcolor=Blue,linkcolor=Blue}
\newcommand{\R}{\mathbb{R}}
\newcommand{\N}{\mathbb{N}}
\newcommand{\E}{\mathbb{E}}
\newcommand{\V}{\mathrm{V}}
\newcommand{\exercice}[3]{\Needspace{6\baselineskip}\vspace{0.5cm}{\color{Blue}\small\textbf{EXERCICE #1}\hfill #2}\par{\large\bfseries #3}\par}

\begin{document}
{\small\color{Blue}PREMIÈRE SPÉCIALITÉ}\par{\LARGE\bfseries Angles et trigonométrie}\par{\large Exercices}\par\vspace{0.2cm}\hrule\vspace{0.4cm}
Les exercices progressent de l’application directe vers la recherche. Détailler les calculs et justifier les réponses.
\par\noindent\begin{minipage}{\linewidth}\exercice{1}{Application}{Conversions}
Convertir $30^\circ$, $225^\circ$ en radians et $7\pi/6$ en degrés.
\end{minipage}\par
\par\noindent\begin{minipage}{\linewidth}\exercice{2}{Application}{Même point}
Ramener $19\pi/6$ et $-7\pi/4$ dans $[0,2\pi[$.
\end{minipage}\par
\par\noindent\begin{minipage}{\linewidth}\exercice{3}{Entraînement}{Valeurs exactes}
Calculer sinus et cosinus de $2\pi/3$ et $7\pi/4$.
\end{minipage}\par
\par\noindent\begin{minipage}{\linewidth}\exercice{4}{Entraînement}{Un signe à choisir}
On sait $\sin x=3/5$ et $x\in[\pi/2,\pi]$. Déterminer $\cos x$.
\end{minipage}\par
\par\noindent\begin{minipage}{\linewidth}\exercice{5}{Entraînement}{Deux solutions}
Résoudre $\cos x=1/2$ dans $[0,2\pi[$.
\end{minipage}\par
\par\noindent\begin{minipage}{\linewidth}\exercice{6}{Application}{Longueur d’arc}
Calculer la longueur d’un arc de rayon $6$ cm correspondant à $\pi/3$ radians.
\end{minipage}\par
\par\noindent\begin{minipage}{\linewidth}\exercice{7}{Approfondissement}{Une affirmation fausse}
La relation $\sin(a+b)=\sin a+\sin b$ est-elle toujours vraie ?
\end{minipage}\par
\par\noindent\begin{minipage}{\linewidth}\exercice{8}{Approfondissement}{Déplacement circulaire}
Un point part de $(1,0)$ et avance de $\pi/3$ radians à chaque étape. Donner ses coordonnées aux étapes $1$, $2$, $3$ et $6$.
\end{minipage}\par
\end{document}
